\section{Threat Model and Security Property}
\label{sec:threatmodel}

\subsection{Attacker Capabilities}
We follow the usual attacker model for TEEs.
We assume a privileged software attacker co-located on the same machine as the victim enclave.
The attacker has compromised the majority of the software stack, including other enclaves and the OS, except for a limited trusted code base (TCB) including the secure bootloader, the SM (together 15K lines of code (LOC)) and the victim program running inside the enclave (see Figure \ref{fig:computing_stack}).
Note that the TCB may contain Spectre gadgets.
The attacker also has control over the contents of shared memory and may control some of the enclave's inputs.
For the rest of this paper, we refer to any code outside the TCB as the OS.
Attackers can exploit the timing of shared microarchitectural events to mount TES attacks.

\subsection{Attacks Considered Out-of-Scope}
\noindent\textbf{Execution-Time Side Channel}
\label{sec:execution_time_SC}
The time an enclave takes to execute might depend on a secret which can be exploited by an attacker.
This includes the timing between any observable (micro)architectural events like syscalls or accesses to shared memory.
We refer to these attacks as exploiting "execution time".
These attacks can only be addressed for constant-time programs \cite{daniel2023prospect, daniel2021hunting, schwarz2020context, choudhary2021speculative, yu2019speculative}. 
To defend against execution-time attacks for this restricted class of programs, Citadel offers the option to disable speculation as we do in the SM (we evaluate a +882\% overhead on a crypto library -- see Section \ref{sec:evalcrypto}).
However, because we mostly want to protect programs that are \textit{not} constant time,
we follow previous work \cite{bourgeat2019mi6, khasawneh2019safespec, kiriansky2018dawg, sakalis2019efficient, li2019conditional, ainsworth2020muontrap, saileshwar2019cleanupspec, wang2019oo7} and consider execution-time attacks out of scope.
These attacks are generally considered lower risk due to lower bandwidth, higher susceptibility to noise, and increased difficulty compared to other side channels. 
For example, NetSpectre \cite{schwarz2019netspectre} can only extract 15 bits per hour using methods applicable to Citadel compared to rates more than 1MB/s for LLC-based side channels \cite{liu2015last}.

\noindent\textbf{Others}
We do not consider physical attacks (e.g., noise, electromagnetic field) or side channels that exploit the physical layer (e.g. power side channels). 
We do not consider Rowhammer attacks.
We are concerned about the privacy of the enclave and do not consider denial-of-service.
As expressed by our contract, we also do not protect enclave programs that explicitly leak their own secrets, including through shared memory.

\subsection{New Building Blocks for Contracts}

\noindent\textbf{Leakage Models}
We define \sw{shm}, a new \emph{weak} leakage model that only represents the trace of accesses to \emph{shared} memory.

\noindent\textbf{Execution Models}
The weakest execution model is \sw{seq}, which represents non-speculative, or sequential execution (one instruction at a time and in-order).
We also have the strongest execution model \sw{spec} which models speculative execution with the following predictors: the pattern history table (PHT), the branch target buffer (BTB) and the return stack buffer (RSB).
We also define \sw{stl} that models straight-line speculation~\cite{straightlinespecARM} where branches are always speculated as non-taken (static branch predictor).
\\

\subsection{Relaxed Microarchitectural Isolation (RMI)}

Given our threat model, we propose a new security property that allow programs to safely share memory.

 \noindent\fbox{%
    \parbox{\dimexpr\columnwidth-2\fboxsep-2\fboxrule\relax}{%
        For a given program, \emph{relaxed microarchitectural isolation} is achieved if information leakage through the microarchitecture is equivalent to the trace of shared memory accesses when running the program non-speculatively (one-instruction-at-a-time and in-order).
    }%
}
That means, for a defense mechanism $\uarch[]{\cdot}{}$, enforcing RMI is exactly equivalent to satisfying $\contract[shm]{\cdot}{seq}$.
Figure \ref{figure:guarantees}, represents how RMI relates to contracts referenced in previous work.
\\

\noindent\textbf{Our goal} is to design microarchitecture that enforces RMI for a large class of programs.
Our design strategy is two-fold.

\begin{enumerate}[leftmargin=*, align=left, itemsep=0ex, topsep=0ex]
    \item Restricting Attacker Observations (HW): we use microarchitectural isolation to restrict the attacker observations to the trace of (speculative) accesses to shared memory.
    \item Restricting Speculative Leakage (HW+SW): we use a combination of mechanisms for controlled speculation and program analysis to ensure these observation traces do not leak more information than if the program was run non-speculatively.
\end{enumerate}

\renewcommand{\ydistance}{0.8}
\renewcommand{\xdistance}{0.75}

\begin{figure}[t!]
    \centering

    \footnotesize
    
    \begin{tikzpicture}[->,>=stealth',shorten >=1pt,auto, semithick]

    % Nodes

    \node[fill=none,draw=none, shape = rectangle, rounded corners, inner sep=0pt, outer sep=0pt, minimum height = 1pt, minimum width = 1pt]  (seqCt) at (0,0) {$\contract[ct]{\cdot}{seq}$};

    \node[fill=none,draw=none, shape = rectangle, rounded corners, inner sep=0pt, outer sep=0pt, minimum height = 1pt, minimum width = 1pt]  (specArch) at ($(seqCt) + (-2*\xdistance,-2.5*\ydistance)$) {$\contract[arch]{\cdot}{spec}$};

    \node[fill=none,draw=none, shape = rectangle, rounded corners, inner sep=0pt, outer sep=0pt, minimum height = 16pt, minimum width = 1pt]  (seqArch) at ($(seqCt) + (-2*\xdistance,-0.5*\ydistance)$) {$\contract[arch]{\cdot}{seq}$};

    \node[fill=none,draw=none, shape = rectangle, rounded corners, inner sep=0pt, outer sep=0pt, minimum height = 16pt, minimum width = 1pt]  (specCt) at ($(seqCt) + (0*\xdistance,-2*\ydistance)$) {$\contract[ct]{\cdot}{spec}$};

     % \node[fill=none,draw=none, shape = rectangle, rounded corners, inner sep=0pt, outer sep=0pt, minimum height = 16pt, minimum width = 1pt]  (seqMemE) at ($(seqCt) + (2*\xdistance,0.5*\ydistance)$) {$\contract[mem]{\cdot}{seq}\sw{+ \epsilon}$};

     % \node[fill=none,draw=none, shape = rectangle, rounded corners, inner sep=0pt, outer sep=0pt, minimum height = 16pt, minimum width = 1pt]  (seqShmE) at ($(seqMemE) + (2*\xdistance,0.5*\ydistance)$) {$\contract[shm]{\cdot}{seq}\sw{+ \epsilon}$};

     \node[fill=none,draw=none, shape = rectangle, rounded corners, inner sep=0pt, outer sep=0pt, minimum height = 16pt, minimum width = 1pt]  (seqMem) at ($(seqCt) + (2*\xdistance,0.5*\ydistance)$) {$\contract[mem]{\cdot}{seq}$};

     \node[fill=none,draw=none, shape = rectangle, rounded corners, inner sep=0pt, outer sep=0pt, minimum height = 16pt, minimum width = 1pt]  (seqShm) at ($(seqMem) + (2*\xdistance,0.5*\ydistance)$) {$\contract[shm]{\cdot}{seq}$};

     % \node[fill=none,draw=none, shape = rectangle, rounded corners, inner sep=0pt, outer sep=0pt, minimum height = 16pt, minimum width = 1pt]  (specShmE) at ($(specCt) + (4*\xdistance,1*\ydistance)$) {$\contract[shm]{\cdot}{spec}\sw{+ \epsilon}$};

     \node[fill=none,draw=none, shape = rectangle, rounded corners, inner sep=0pt, outer sep=0pt, minimum height = 16pt, minimum width = 1pt]  (specShm) at ($(specCt) + (4*\xdistance,1*\ydistance)$) {$\contract[shm]{\cdot}{spec}$};

      \node[fill=none,draw=none, shape = rectangle, rounded corners, inner sep=0pt, outer sep=0pt, minimum height = 16pt, minimum width = 1pt]  (stlShm) at ($(specShm) + (0*\xdistance,1*\ydistance)$) {$\contract[shm]{\cdot}{stl}$};

      % \node[fill=none,draw=none, shape = rectangle, rounded corners, inner sep=0pt, outer sep=0pt, minimum height = 16pt, minimum width = 1pt]  (botE) at ($(specShmshiE) + (2*\xdistance,-0.5*\ydistance)$) {$\contract[\bot]{\cdot}{}\sw{+ \epsilon}$};

      \node[fill=none,draw=none, shape = rectangle, rounded corners, inner sep=0pt, outer sep=0pt, minimum height = 16pt, minimum width = 1pt]  (bot) at ($(seqShm) + (2*\xdistance,0*\ydistance)$) {$\contract[\bot]{\cdot}{}$};

    % Arrows

    \path (seqArch) edge[] node[left] {} (seqCt);
    \path (specArch) edge[] node[left] {} (seqArch);
    \path (specArch) edge[] node[left] {} (specCt);
    \path (specCt) edge[] node[left] {} (seqCt);
    \path (seqCt) edge[] node[left] {} (seqMem);
    \path (specCt) edge[] node[left] {} (specShm);
    \path (seqMem) edge[] node[left] {} (seqShm);
    % \path (specShmE) edge[] node[left] {} (seqShmE);
    % \path (specShmE) edge[] node[left] {} (specShm);
    % \path (specShmE) edge[] node[left] {} (botE);
    % \path (botE) edge[] node[left] {} (bot);
    \path (seqShm) edge[] node[left] {} (bot);
    \path (specShm) edge[] node[left] {} (stlShm);
    \path (stlShm) edge[] node[left] {} (seqShm);

    % uArch
    
    \begin{scope}[on background layer]
        % \node[draw=gray!50,dashed, fill=none, fill opacity=0.5, rounded corners, rotate fit=-8
        % , minimum height=0.1em, fit=(seqArch)](FitLd) {};
        % \node  (label1) at ($(FitLd.east) + (0.5,0)$) {$\LoadDelayMuarchSem$};

        \node[draw=gray!50,dashed, fill=none, fill opacity=0.5, rounded corners, rotate fit=0, minimum height=0.1em, fit=(seqCt)](FitCt) {};
        \node  (label1) at ($(FitCt.north) + (0,0.25)$) {$\uarchf[pro]{\cdot}{}$};
        
        \node[draw=gray!50,dashed, fill=none, fill opacity=0.5, rounded corners, rotate fit=-36, minimum height=0.1em, fit=(seqArch) (specCt)](FitTt) {};
        \node  (label2) at ($(FitTt.east) + (0.5,0)$) {$\uarchf[tt]{\cdot}{}$};
        
        % \node[draw=gray!50,dashed, fill=none, fill opacity=0.5, rounded corners, rotate fit=+14, minimum height=0.1em, fit=(seqCt) (seqArch) (seqMemE) (seqShmE) (seqShm)](FitSeq) {};
        % \node  (label3) at ($(FitSeq.east) + (0.7,0)$) {$\uarchf[seq]{\cdot}{}$};

        % \node[draw=gray!50,dashed, fill=none, fill opacity=0.5, rounded corners, rotate fit=0, minimum height=0.1em, fit=(seqShmE)](FitSafe*) {};
        % \node  (label4) at ($(FitSafe*.north) + (0,0.25)$) {$\uarchf[safe*]{\cdot}{}$};

        \node[draw=red!100,dashed, fill=none, fill opacity=0.5, rounded corners, rotate fit=0, minimum height=0.1em, fit=(seqShm)](FitSafe) {};
        % \node  (label5) at ($(FitSafe.north) + (0,0.25)$) {$\uarch[safe]{\cdot}{}$};
        % \node  (label6) at ($(FitSafe.east) + (0.7,0)$) {$\uarch[burst]{\cdot}{sta}$};
        \node  (label5) at ($(FitSafe.north) + (0,0.25)$) {$\uarch[safe]{\cdot}{} \hw{/} \uarch[burst]{\cdot}{sta}$};

        \node[draw=gray!50,dashed, fill=none, fill opacity=0.5, rounded corners, rotate fit=0, minimum height=0.1em, fit=(stlShm)](FitBurstRaw) {};
        \node  (label7) at ($(FitBurstRaw.east) + (0.5,0)$) {$\uarchf[burst]{\cdot}{\top}$};

        % \node[draw=gray!50,dashed, fill=none, fill opacity=0.5, rounded corners, rotate fit=0, minimum height=0.1em, fit=(botE)](FitMI6*) {};
        % \node  (label8) at ($(FitMI6*.south) + (0,-0.25)$) {$\uarchf[mi6*]{\cdot}{}$};

        \node[draw=gray!50,dashed, fill=none, fill opacity=0.5, rounded corners, rotate fit=0, minimum height=0.1em, fit=(bot)](FitMI6) {};
        \node  (label9) at ($(FitMI6.north) + (0,0.25)$) {$\uarchf[mi6]{\cdot}{}$};
    \end{scope}


    \end{tikzpicture}
    \caption{Security guarantees of secure-speculation mechanisms. 
    Hardware semantics of Citadel' defense modes, Safe ($\uarch[safe]{\cdot}{}$) and Burst ($\uarch[burst]{\cdot}{sta}$), both satisfy RMI ($\contract[shm]{\cdot}{seq}$).
    We also represented other relevant hardware semantics such as $\uarch[burst]{\cdot}{\top}$ which represents the hardware mechanisms of Burst without the associated static code analysis (see Section \ref{sec:burstsecanalysis}), and some previous works such as $\uarch[pro]{\cdot}{}$ for ProSpeCT\cite{daniel2023prospect} or $\uarch[tt]{\cdot}{}$ for taint-tracking sandboxing schemes such as STT~\cite{yu2019speculative}. %
    }\label{figure:guarantees}
\end{figure}
